\chapter[APRESENTAÇÃO DA TRAJETÓRIA DO ALUNO]{APRESENTAÇÃO DA TRAJETÓRIA DO ALUNO}

Em 2017, iniciei minha jornada na Pontifícia Universidade Católica do Rio Grande do Sul (\ac{pucrs}), no curso de Engenharia de Computação. Foi nesse momento que tive meu primeiro contato com programação e me encantei com as possibilidades que ela proporcionava. Em 2018, durante um evento de Portas Abertas da \ac{pucrs}, conheci o curso de Engenharia de Software e percebi que ele abrangia muito mais a parte de programação do que a Engenharia de Computação. Decidi solicitar a troca de curso em 2019.

No segundo semestre de 2019, tentei cursar a disciplina de \ac{ages} I pela primeira vez. No entanto, logo no primeiro mês, acabei trancando a disciplina, pois senti que não estava conseguindo contribuir e enfrentei algumas questões pessoais. Tive minha segunda oportunidade como \ac{ages} I em 2023/01 e senti-me muito mais preparado para enfrentá-la do que antes. Nesse momento, já possuía conhecimentos mais aprofundados e pude auxiliar ativamente a equipe.

Em 2024/01, durante minha segunda experiência com a \ac{ages}, pude contribuir ainda mais para o projeto, sem a necessidade de suporte constante de colegas de semestres superiores. Meu crescimento me permitiu assumir responsabilidades em tarefas mais complexas, contribuindo para acelerar o processo de desenvolvimento e marcando minha evolução como estudante de Engenharia de Software.

Em 2025/01, assumi o papel de \ac{ages} III. Nesta posição, pude aplicar todo o conhecimento acumulado nos semestres anteriores, agora como mentor e líder técnico. A responsabilidade de orientar \ac{ages} I e II, tomar decisões arquiteturais e garantir a qualidade das entregas consolidou minha transformação de aprendiz para orientador. A capacidade de resolver problemas críticos sob pressão demonstraram que estava pronto para os desafios do mercado de trabalho.

Em 2026/01, assumi o papel de \ac{ages} IV no projeto Estudify. Foi a primeira vez que olhei para um projeto inteiramente pela ótica da gestão, e não mais pela da execução. Em vez de pegar tarefas para desenvolver, meu trabalho passou a ser escrever e detalhar as histórias, distribuí-las entre os integrantes da minha \textit{squad}, acompanhar o andamento das entregas e remover impedimentos quando apareciam. Aprender a confiar no time e a delegar as tarefas, ao invés de tentar executá-las eu mesmo, foi um dos maiores desafios dessa experiência. Esse semestre fechou o ciclo da \ac{ages} para mim, consolidando uma trajetória que começou com dúvidas sobre onde eu poderia contribuir e terminou com a responsabilidade de garantir que um time inteiro chegasse junto nas entregas.
